% Prijava teme projekta — RSO 2026/27, Projektna ekipa 13. Izris: make (lualatex)
\documentclass[11pt,a4paper]{article}
\usepackage[top=2.3cm,bottom=2.3cm,left=2.4cm,right=2.4cm,headsep=0.8cm]{geometry}
\usepackage{fontspec}
\setmainfont{Libertinus Serif}[Extension=.otf, UprightFont=LibertinusSerif-Regular, BoldFont=LibertinusSerif-Semibold,
  ItalicFont=LibertinusSerif-Italic, BoldItalicFont=LibertinusSerif-SemiboldItalic]
\setsansfont{TeX Gyre Heros}[Extension=.otf, UprightFont=texgyreheros-regular, BoldFont=texgyreheros-bold,
  ItalicFont=texgyreheros-italic, BoldItalicFont=texgyreheros-bolditalic, Scale=0.92]
\usepackage{polyglossia}
\setdefaultlanguage{slovenian}
\usepackage{xcolor}
\definecolor{ink}{HTML}{1F2A37}
\definecolor{accent}{HTML}{2F5E9E}
\definecolor{muted}{HTML}{6B7280}
\definecolor{rule}{HTML}{D1D5DB}
\usepackage{graphicx,booktabs,tabularx,array,enumitem,caption}
\usepackage[hidelinks]{hyperref}
\urlstyle{sf}
\usepackage{titlesec}
\usepackage{fancyhdr}
\usepackage[document]{ragged2e}

\titleformat{\section}{\sffamily\large\bfseries\color{ink}}{\color{accent}\thesection}{0.7em}{}
\titleformat{\subsection}{\sffamily\normalsize\bfseries\color{ink}}{}{0em}{}
\titlespacing*{\section}{0pt}{1.3em}{0.5em}
\titlespacing*{\subsection}{0pt}{0.8em}{0.3em}
\setlist{nosep, leftmargin=1.4em, itemsep=1pt}
\setlength{\parindent}{0pt}
\setlength{\parskip}{5pt}
\setlength{\emergencystretch}{3em}
\captionsetup{font={small,sf}, labelfont={bf,color=accent}, labelsep=period, skip=6pt}
\renewcommand{\arraystretch}{1.25}
\newcolumntype{Y}{>{\RaggedRight\arraybackslash}X}

\pagestyle{fancy}
\fancyhf{}
\fancyhead[L]{\sffamily\footnotesize\color{muted} Šepet — prijava teme projekta}
\fancyhead[R]{\sffamily\footnotesize\color{muted} RSO 2026/27 · Projektna skupina 13}
\fancyfoot[C]{\sffamily\footnotesize\color{muted}\thepage}
\renewcommand{\headrulewidth}{0pt}
\fancypagestyle{first}{\fancyhf{}\fancyfoot[C]{\sffamily\footnotesize\color{muted}\thepage}}

\begin{document}
\thispagestyle{first}

{\sffamily\color{muted}\small Prijava teme projekta · Računalniške storitve v oblaku 2026/27 · UL FRI}\par\vspace{6pt}
{\sffamily\bfseries\fontsize{24}{28}\selectfont\color{ink} Šepet}\par\vspace{2pt}
{\sffamily\Large\color{ink} Varno sporočanje s šifriranjem od konca do konca}\par\vspace{10pt}
{\color{rule}\rule{\textwidth}{0.6pt}}\par\vspace{2pt}
{\small
\begin{tabularx}{\textwidth}{@{}>{\sffamily\color{muted}}l Y@{}}
Skupina & Projektna skupina 13 — Ključarji \\
Člani & Tibor Koderman, Amar Mehić, Eva Vidic \\
Repozitoriji & \url{https://gitlab.koderverse.com/fri-mag/rso} (en repozitorij za vsako mikrostoritev) \\
\end{tabularx}}\par
{\color{rule}\rule{\textwidth}{0.6pt}}

\section{Opis projekta}
Pri večini aplikacij za klepet moramo ponudniku zaupati, da sporočil ne bo bral ali jih predal tretjim osebam; Telegram na primer klepetov privzeto ne šifrira od konca do konca, zahteve po dostopu oblasti do zasebnih sporočil pa so vse pogostejše. Razvili bomo aplikacijo Šepet za sporočanje ena na ena in v skupinah z oblikovanim besedilom in slikami. Vsebina se šifrira na napravi pošiljatelja s standardnim protokolom MLS (RFC 9420) in dešifrira šele na napravah prejemnikov, zato strežnik hrani in posreduje le šifrirane podatke. Sistem bomo zgradili iz sedmih mikrostoritev in ga namestili v oblak Azure (Kubernetes).

\section{Varnostni model}
Strežniku zaupamo dostavo, ne pa zaupnosti:
\begin{itemize}
\item \textbf{Vse v pogovoru je šifrirano} — sporočila, slike in imena skupin, pa tudi tipkanje, potrdila o branju in reakcije. Strežnik vidi le, kdo je v katerem pogovoru.
\item \textbf{Pregleden dnevnik ključev} — javni ključi naprav so v dnevniku z Merklovim drevesom, zato strežnik zamenjave ključa ne more prikriti; uporabnika lahko primerjata tudi varnostno številko s kodo QR.
\item \textbf{Prijava brez gesel} — s ključi passkey (WebAuthn).
\item \textbf{Obvestila brez vsebine} — potisno obvestilo napravo le zbudi; sporočilo prenese, dešifrira in prikaže aplikacija sama.
\item \textbf{Izginjajoča sporočila} — po izteku časovnika jih izbrišejo naprave in strežnik.
\end{itemize}

\section{Tehnologije in razvojno okolje}
{\small
\begin{tabularx}{\textwidth}{@{}>{\sffamily}l Y@{}}
\toprule
Odjemalec & Flutter (Android in splet) \\
Šifriranje & MLS (OpenMLS), rezervna možnost Signal Protocol \\
Mikrostoritve & .NET 10 (ASP.NET Core), Clean Architecture \\
Podatkovne baze & PostgreSQL, ScyllaDB, Azure Blob Storage \\
Komunikacija & REST, gRPC (Protobuf), WebSocket, NATS JetStream; GraphQL samo za skrbniško stran \\
Identiteta & lasten ponudnik OAuth 2.0 / OpenID Connect v .NET (OpenIddict, passkeys) \\
Platforma & Docker, Kubernetes (Azure AKS), Helm, Envoy Gateway, Terraform \\
CI/CD in opazovanje & GitLab CI; Prometheus, Grafana, Loki \\
Drugo & Azure Functions, Firebase Cloud Messaging \\
Orodja & JetBrains Rider, VS Code, Android Studio, Docker Compose, kind, Git \\
\bottomrule
\end{tabularx}}

\section{Arhitektura}
Aplikacija dostopa do storitev prek prehoda API, ki preverja žetone JWT in usmerja zahteve (REST, WebSocket). Kadar storitev potrebuje takojšen odgovor, kliče drugo storitev prek gRPC (npr. preverjanje članstva pri pošiljanju sporočila); dejstva, ki zanimajo več storitev, objavi kot dogodek v NATS JetStream. Vsaka storitev s trajnimi podatki ima svojo bazo, do podatkov drugih storitev pa pride le prek API-jev ali lokalnih kopij iz dogodkov; šifrirane slike potujejo neposredno v Azure Blob Storage (slika~\ref{fig:arhitektura}). Pogodbe so v repozitoriju vsake storitve: OpenAPI za REST, Protobuf za gRPC in AsyncAPI za dogodke.

\begin{figure}[ht]
\centering
\includegraphics[width=\textwidth]{build/arhitektura.pdf}
\caption{Začetna shema arhitekture s komunikacijskimi protokoli.}
\label{fig:arhitektura}
\end{figure}

\section{Mikrostoritve}
{\small
\begin{tabularx}{\textwidth}{@{}>{\sffamily}l Y >{\RaggedRight}p{2.1cm} l@{}}
\toprule
\textbf{Storitev} & \textbf{Funkcionalnosti} & \textbf{Baza} & \textbf{Odgovorna oseba} \\
\midrule
identity & ponudnik OIDC: prijava s passkey, organizacije, vloge, naprave; skrbniški GraphQL (branje) & PostgreSQL & Tibor Koderman \\
keys & javni ključi naprav, paketi ključev MLS, pregleden dnevnik ključev & PostgreSQL & Tibor Koderman \\
conversation & pogovori in skupine, člani, prošnje za pogovor, časovnik izginjanja; Event Sourcing in CQRS & PostgreSQL & Amar Mehić \\
message & šifrirana sporočila z vrstnim redom po epohi MLS, sinhronizacija naprav, samodejno brisanje & ScyllaDB & Amar Mehić \\
realtime & WebSocket: dostava v živo, prisotnost, šifrirani efemerni dogodki & — & Amar Mehić \\
media & nalaganje in prenos šifriranih slik, omejitve velikosti & PostgreSQL & Eva Vidic \\
notification & potisna obvestila brez vsebine (Firebase Cloud Messaging) & PostgreSQL & Eva Vidic \\
\bottomrule
\end{tabularx}}

Poleg tega bomo razvili aplikacijo v Flutterju in funkcijo Azure Functions, ki preveri naložene slike in briše potekle (Eva Vidic), ter skupno platformo s prehodom API, CI/CD in opazovanjem (Tibor Koderman).

\section{Skaliranje}
Mikrostoritve ne hranijo stanja, zato jih Kubernetes lahko poljubno podvaja, ko pride več uporabnikov:
\begin{itemize}
\item \textbf{HorizontalPodAutoscaler} doda pode storitvam identity, conversation, message in media, ko naraste obremenitev procesorja.
\item \textbf{KEDA} skalira realtime po številu odprtih povezav WebSocket in notification po zaostanku dogodkov v NATS.
\item \textbf{Samodejno skaliranje vozlišč AKS} doda vozlišče, ko novi podi nimajo prostora.
\item Realtime prejema sporočila prek NATS, zato je uporabnik lahko povezan s katerimkoli podom; ScyllaDB se skalira z dodajanjem vozlišč.
\end{itemize}
Skaliranje bomo pokazali z obremenitvenim testom (k6) in metrikami v Grafani.

\section{Namestitev in stroški}
\begin{itemize}
\item \textbf{Produkcija} teče v eni majhni gruči AKS (2 vozlišči z 2 vCPU), ki jo ustavimo, ko je ne uporabljamo. Baze tečejo v gruči (CloudNativePG, ScyllaDB), slike vsebnikov so v registru GitLab, zato ne plačujemo upravljanih baz, registra ali Azure Monitor. Ocena je 15–25 USD na mesec z opozorilom proračuna pri 25 USD, kar je za ves semester manj kot kredit enega študenta.
\item \textbf{Razvoj} poteka lokalno (Docker Compose, kind); cevovodi CI tečejo na našem strežniku GitLab, dodaten izvajalec CI pa na domačem strežniku. Vse ocenjevane storitve ostanejo v Azure.
\end{itemize}

\section{Primeri uporabe}
\begin{itemize}
\item registracija in prijava s ključem passkey;
\item iskanje uporabnika in prošnja za pogovor, ki jo prejemnik sprejme ali zavrne;
\item pošiljanje sporočil z oblikovanim besedilom in stanjem dostave (poslano, dostavljeno, prebrano);
\item pošiljanje slik;
\item skupine: ustvarjanje, dodajanje in odstranjevanje članov;
\item obvestila, ko aplikacija ni odprta;
\item preverjanje varnostne številke s kodo QR;
\item izginjajoča sporočila;
\item upravljanje organizacije in vlog uporabnikov (skrbnik).
\end{itemize}

\subsection{Kompleksnejši primer: slika v skupinskem pogovoru}
\begin{enumerate}
\item Aplikacija sliko šifrira z naključnim ključem in pri storitvi \textsf{media} dobi povezavo za nalaganje.
\item Slika se naloži v Azure Blob Storage; funkcija Azure jo preveri in obvesti \textsf{media}.
\item Aplikacija pošlje šifrirano sporočilo s ključem slike storitvi \textsf{message}, ta prek gRPC pri \textsf{conversation} preveri članstvo pošiljatelja.
\item \textsf{message} sporočilo shrani v ScyllaDB in objavi dogodek v NATS JetStream.
\item \textsf{realtime} ga dostavi povezanim članom, \textsf{notification} pa ostalim pošlje potisno obvestilo.
\item Prejemniki sporočilo in sliko dešifrirajo na svojih napravah.
\end{enumerate}

\end{document}
